\documentclass[12pt,a4paper]{article}

\usepackage[utf8]{inputenc}
\usepackage[T1]{fontenc}
\usepackage{lmodern}
\usepackage{geometry}
\usepackage{amsmath}
\usepackage{parskip}
\usepackage{mathbbol}
\usepackage{graphicx}
\usepackage{pdfpages}
\usepackage{setspace}
\geometry{margin=1in}

\title{}
\author{}
\date{}

\begin{document}
\doublespacing
\maketitle

\section{Optimal Control Problem}

Based on \texttt{mpc\_landing/mpc.py}:

\begin{equation}
\begin{aligned}
\min_{\{x[t],\,u[t]\}_{t=0}^{N-1}} \quad 
& \sum_{t=0}^{N-1} \Big[
(x[t]-x_{\mathrm{ref}}[t])^\top Q (x[t]-x_{\mathrm{ref}}[t])
+ (u[t]-u_{\mathrm{eq}})^\top R (u[t]-u_{\mathrm{eq}})
\Big] \\
& \quad + (x[N]-x_{\mathrm{ref}}[N])^\top Q_f (x[N]-x_{\mathrm{ref}}[N]) \\
\text{s.t.} \quad
& x[t+1] = A x[t] + B u[t] + g, \quad t = 0,\dots,N-1, \\
& x[0] = x_0^{\mathrm{aug}}, \\
& -a_{\max}\mathbf{1} \le u[t] \le a_{\max}\mathbf{1}, \quad t = 0,\dots,N-1, \\
& -v_{\max} \le x_{2i+1}[t] \le v_{\max}, \quad i=0,1,2,\; t = 1,\dots,N.
\end{aligned}
\end{equation}

\begin{equation}
x[t] =
\begin{bmatrix}
p_x[t] \\ v_x[t] \\ p_y[t] \\ v_y[t] \\ p_z[t] \\ v_z[t] \\ d[t]
\end{bmatrix},
\qquad
u[t] =
\begin{bmatrix}
a_x[t] \\ a_y[t] \\ a_z[t]
\end{bmatrix}.
\end{equation}

\begin{equation}
u_{\mathrm{eq}} =
\begin{bmatrix}
0 \\ g \\ 0
\end{bmatrix}.
\end{equation}

In this optimal control problem, the state vector $x[t]\in\mathbb{R}^7$ represents the system configuration at time step $t$, consisting of positions and velocities in three spatial directions together with a disturbance term, i.e.,
\[
x[t]=[p_x[t],\,v_x[t],\,p_y[t],\,v_y[t],\,p_z[t],\,v_z[t],\,d[t]]^\top,
\]
where $d[t]$ models a constant bias in vertical acceleration. The control input $u[t]\in\mathbb{R}^3$ denotes the commanded accelerations along each axis,
\[
u[t]=[a_x[t],\,a_y[t],\,a_z[t]]^\top.
\]
The matrices $A$ and $B$ define the discrete-time linear dynamics of a double-integrator system, while the vector $g$ accounts for the effect of gravity as an affine term. The initial condition $x_0^{\mathrm{aug}}$ is the current measured state augmented with the disturbance estimate. The reference trajectory $x_{\mathrm{ref}}[t]$ specifies the desired positions and velocities over the prediction horizon, and the objective function penalizes deviations from this reference using the positive semidefinite weighting matrices $Q$ and $Q_f$ for stage and terminal costs, respectively. The matrix $R$ penalizes control effort relative to the equilibrium input $u_{\mathrm{eq}}$, which corresponds to the hover condition that balances gravity. The scalar bounds $a_{\max}$ and $v_{\max}$ impose infinity-norm constraints on the control inputs and velocities, ensuring that accelerations and speeds remain within physically feasible limits. The horizon length $N$ defines the number of discrete time steps over which the optimization is performed.



\end{document}
